\chapter{Introduction}
\label{ch:introduction}

\section{Motivation}
\label{sec:intro-motivation}

Previous research studied the development of shared abstractions in a collaborative two-player game called the architect-builder game, in which one player (the architect) sees a target scene and instructs the other (the builder), who cannot see it, on how to reconstruct it. McCarthy et al.~\cite{mccarthy_learning_2021} showed that human pairs playing this game develop local shortcuts for recurring structures, while Thomas et al.~\cite{thomas_pace_2025} showed with PACE that analogous abstractions can emerge in neural agents without prior visual or linguistic knowledge. Chapter~\ref{ch:background} gives the full background, here the important point is that PACE provides a bias-free setting in which the formation of reusable communicative units can be inspected directly.

This thesis extends these results in two directions. First, natural input is not uniform. If you take a random English text and count how often each word appears, you find that a very small number of words (``the'', ``of'', ``and'') account for most of the occurrences, while the vast majority of words appear only a handful of times. This heavily skewed pattern is called a Zipfian distribution, and it is one of the most robust regularities in natural language~\cite{zipf_human_1949}. Frequent words also tend to be shorter. That length is driven less by frequency itself than by a word's predictability in context~\cite{piantadosi_word_2011}. Word length depends less on how common a word is than on how easy it is to guess from context: an easy-to-guess word carries little new information, so it can be short. A word like ``the'' is a good example: it is short not just because it is common, but because it is easy to predict, since a noun almost always comes next. A lesser common word is ``molybdenum'', it is not easy to predict in a phrase. Both effects are usually read as signatures of a pressure that balances informativeness against production cost~\cite{gibson_how_2019}. A uniform training distribution removes precisely this coupling of frequency, predictability, and form. Second, different communities do not share exactly the same distribution of experience. We therefore test PACE under distributions of different concentrations and under cross-population pairings, characterising how input concentration shapes the vocabularies agents build and how robust those vocabularies are when agents trained under different conditions meet.

\section{Contributions and Research Questions}
\label{sec:intro-contributions}

This thesis extends PACE~\cite{thomas_pace_2025} in two directions. On the input side, it replaces PACE's uniform training distribution with a family of Zipfian distributions, parameterised by a concentration parameter $\alpha$ that lets us move smoothly from the original uniform case ($\alpha = 0$) to a highly concentrated one ($\alpha = 3$). On the analysis side, it develops three research questions, each corresponding to a different level at which the input distribution might shape emergent communication: at the level of the language a single population develops (RQ1), at the level of the linguistic differences between populations (RQ2), and at the level of the communicative interactions between populations (RQ3).

\paragraph{RQ1: the internal structure of a single population's language.}
The first research question asks how the concentration $\alpha$ shapes the emergent language a single population develops. We look at three things: how many abstractions the bandit algorithm adopts into its vocabulary, how the language compresses over training, and what signal (raw frequency, or the MDL score that trades frequency against length) drives the bandit's adoption decisions. Our expectation is that a more concentrated distribution should give the bandit fewer distinct high-frequency patterns to work with, and therefore push it toward a smaller and more specialised library. Chapter~\ref{ch:experimental} tests this hypothesis and finds a pattern more nuanced than a simple monotonic decrease: three regimes emerge across the $\alpha$ range, with the vocabulary being largest and most stable in the intermediate regime and impoverished at both ends.

\paragraph{RQ2: dialect formation across populations.}
The second research question asks whether populations that share the concentration $\alpha$ but differ in the ranking permutation of the Zipfian dataset develop different abstraction libraries. Because the ranking permutation determines which specific scenes end up being most frequent in each population's training data, two populations with the same $\alpha$ but different permutations see a distribution with the same shape but with different scenes at each position of the ranking. Any similarity in the resulting libraries therefore has to come from structural convergence rather than from identical training data. We test this by training twenty PACE populations at each $\alpha$, measuring the pairwise overlap between their libraries, and comparing that overlap to a null model that captures the baseline similarity we would expect by chance. Chapter~\ref{ch:rq2} shows that the overlap decreases monotonically with $\alpha$: at low $\alpha$, populations largely agree on their vocabularies; at high $\alpha$, they diverge. Qualitative analysis at moderate $\alpha$ traces the divergence back to the specific pairs that happen to be frequent in each population's dataset.

\paragraph{RQ3: cross-population communication.}
The third research question asks whether, and if so how, these library-level differences shape the adaptation dynamics when two different populations are paired together. If the languages developed by different populations are really local to their training data, we should expect cross-population communication to be harder than within-population, and the difficulty should scale with the distance between the two training distributions. Chapter~\ref{ch:rq3} tests this along three axes: pairings between agents from the same $\alpha$ but different rankings (Group C), pairings between agents from different $\alpha$ values (Group X), and case studies further removed from the training conditions such as mixed distributions and out-of-distribution shifts (Group O). Along the way, three modifications to the same-$\alpha$ setup (Group M) let us look inside the adaptation dynamics and see what the bandit is actually doing during joint training. The results give a picture that differs from what we expected: adaptation difficulty in the same-$\alpha$ setting depends strongly on the shared concentration regime and only weakly on the pairwise distance between the two agents, under any of the three distance metrics we tested. When the two agents come from different $\alpha$ values, the difficulty is asymmetric and driven by the specificity of the speaker's abstraction library rather than by the neural mismatch between the two agents. The picture that emerges is that cross-population robustness is a structural property of the training conditions, and that neural adaptation alone cannot compensate when the two agents' discrete vocabularies are too poorly overlapping to support communication.

\section{Thesis Structure}
\label{sec:intro-thesis-structure}

Chapter~\ref{ch:background} reviews the background on emergent communication research and situates PACE in that landscape. Chapter~\ref{ch:pace} describes PACE in detail and identifies the specific components that this thesis modifies. Chapter~\ref{ch:technical} describes the experimental framework: the datasets, the metrics, the training procedure, and the computational setup used to run the experiments at scale.

The three empirical chapters follow. Chapter~\ref{ch:experimental} tests RQ1: it characterises how the concentration $\alpha$ affects the internal structure of the emergent language, from adoption counts to compression dynamics. Chapter~\ref{ch:rq2} tests RQ2: it measures dialect formation by comparing the abstraction libraries of twenty populations at each $\alpha$, and analyses the qualitative reasons for the observed divergence. Chapter~\ref{ch:rq3} tests RQ3 through the four experimental groups introduced above: same-$\alpha$ cross-pairing, modifications to the base setup that expose the mechanics of adaptation, cross-$\alpha$ pairing, and case studies of setups further from the training conditions.

Chapter~\ref{ch:conclusion} draws together the three empirical chapters and returns to the broader question that motivates the work: what emergent communication research can tell us about the conditions under which shared linguistic conventions form, and about the conditions under which they break.